\documentclass[12pt]{scrartcl}
\usepackage{fontawesome5}
\usepackage{siunitx}

%%%%%%%%%%%%%%%%%%%%%%%%%%%%%%%%%%
%        Graphics & boxes        %
%%%%%%%%%%%%%%%%%%%%%%%%%%%%%%%%%%

\usepackage{xcolor}
\usepackage{hyperref}
\usepackage{catppuccinpalette}
\usepackage{tikz}
\usepackage[breakable]{tcolorbox}
\tcbuselibrary{theorems, skins}

\tikzset{
	second box/.style={
			anchor=east,
			text=white,
			% rounded corners,
			fill=#1,
			xshift=-4mm,
		},
	tcbreak/.style = {fill=white, thick, draw=#1, text=#1, font=\bfseries}
}

\tcbset{
	common/.style n args={2}{
            parbox=false,
			colframe={#1},
			colback={#1!5},
			colbacktitle={#1},
			lower separated=false,
			coltitle=white,
			boxrule=1.25pt,
			fonttitle=\bfseries,
			enhanced,
			breakable,
			top=10pt,
			sharp corners,
			before skip=1em,
			after skip=2em,
			attach boxed title to top left={
					yshift=-0.25cm,
					xshift=0.38cm,
				},
			boxed title style={
					boxrule=0pt,
					colframe=white,
					arc=0pt,
					outer arc=0pt,
				},
			separator sign={~~},
			% -- overlays based on the following TeX SE answer: https://tex.stackexchange.com/a/545324/162087
			overlay unbroken={
					\node[text=white, align=right, fill=#1, xshift=-4mm, minimum height=6mm, anchor=east] at (frame.south east) {#2};
				},
			overlay first={
					\draw[thick, #1, decoration={coil, amplitude=0.5mm}, decorate]
					(frame.south west) -- (frame.south east);
					\path[font=\small\itshape] (frame.south) node [tcbreak={#1}] (cont) {continues in the next page \faHandPointRight};
				},
			overlay middle={%
					% upper break
					\draw[thick, #1, decoration={coil, amplitude=0.5mm}, decorate]
					(frame.north west) -- (frame.north east);
					\path[font=\small\itshape] (frame.north) node [tcbreak={#1}] (cont) {\faHandPointRight continues from the previous page};

					% lower break
					\draw[thick, #1, decoration={coil, amplitude=0.5mm}, decorate]
					(frame.south west) -- (frame.south east);
					\path[font=\small\itshape] (frame.south) node [tcbreak={#1}] (cont) {continues in the next page \faHandPointRight};
				},
			overlay last={
					\node[text=white, align=right, fill=#1, xshift=-4mm, minimum height=6mm, anchor=east] at (frame.south east) {#2};
					\draw[thick, #1, decoration={coil, amplitude=0.5mm}, decorate]
					(frame.north west) -- (frame.north east);
					\path[font=\small\itshape] (frame.north) node [tcbreak={#1}] (cont) {\faHandPointRight continues from the previous page};
				}
		},
	examplestyle/.style={
			common={CtpLatteBlue}{\faStar},
		},
}

\newcommand{\newterm}[1]{\textcolor{CtpLatteMauve}{\textbf{#1}}}

% args are {environment name}{Title}{style}{reference prefix}
\newtcbtheorem[auto counter, number within=section]{example}{Example}{examplestyle}{example}
%
%%%%%%%%%%%%%%%%%%%%%%%%
%        Minted        %
%%%%%%%%%%%%%%%%%%%%%%%%

\usepackage{minted}
\setminted[python]{
    bgcolor=CtpLatteBase,
    % frame=lines,
    linenos,
}
% \usemintedstyle{}
\newcommand{\pyinl}[1]{\mintinline{python}{#1}}

\setlength{\parindent}{0pt}
\setlength{\parskip}{1em}


%%%%%%%%%%%%%%%%%%%%%%%%%%%%%%%
%        Maths related        %
%%%%%%%%%%%%%%%%%%%%%%%%%%%%%%%

\usepackage{amsmath, amsfonts, amssymb}
\newcommand{\mtrx}[1]{\mathbf{#1}}
\newcommand{\vvec}[1]{\vec{\mathrm{#1}}}
\newcommand{\uvec}[1]{\hat{\mathrm{#1}}}
\newcommand{\T}[1]{#1^{\intercal}}

\colorlet{xcomp}{CtpLatteRed}
\colorlet{ycomp}{CtpLatteBlue}
\colorlet{zcomp}{CtpLatteGreen}

% Row- and column-vectors: arguments separated by ";".
% Example: $\vec{a} = \colvec{1;2;3;4}$.
\makeatletter
\newcommand\rcvector[2][\\]{\ensuremath{%
		\global\def\rc@delim{#1}%
		\negthinspace\begin{bmatrix}
			\rc@vector #2;\relax\noexpand\@eolst%
		\end{bmatrix}}}
\def\rc@vector #1;#2\@eolst{%
	\ifx\relax#2\relax
        #1
	\else
        #1\rc@delim
		\rc@vector #2\@eolst%
	\fi}
\makeatother
\newcommand{\colvecxy}[2]{%
    \begin{bmatrix}
        \textcolor{xcomp}{#1}\\
        \textcolor{ycomp}{#2}
    \end{bmatrix}
}
\newcommand{\rowvec}[1]{\rcvector[\ \;]{#1}}

\newcommand{\PC}[1]{\mathrm{PC}_{#1}}
\newcommand{\var}[1]{\mathrm{Var}_{#1}}


%%%%%%%%%%%%%%%%%%%%%%%%%%
%        Doc data        %
%%%%%%%%%%%%%%%%%%%%%%%%%%

\title{Softwaregestützte Modellbildung und Simulation: Exercise 2}
\subtitle{Principal Component Analysis (PCA)}
\author{Peleg Bar Sapir}

%%%%%%%%%%%%%%%%%%%%%%%%%%%%%%%
%        Document here        %
%%%%%%%%%%%%%%%%%%%%%%%%%%%%%%%

\begin{document}
\maketitle

Principal Component Analysis (PCA) is an application of Singular Value Decomposition (SVD) to a data set, which allows us to visualize the data by \textit{dimensionality reduction}: instead of looking at a high-dimensional data (sometimes in excess of several \textit{hundreds} dimensions!), which may have numerous correlations - we can look at projections of said data on the first few directions in which its variance is the largest.

In standard PCA, the data set is represented as a matrix, with its \textbf{rows} being the separate \textit{samples} (measurements), and its \textbf{columns} the different \textit{variables} (features).

The following steps describe how to implement PCA in Python using NumPy's \pyinl{SVD()} function, and visualize the resulting data representation using MatPlotLib's \pyinl{axes.bars()}, \pyinl{axes.plot()} and \pyinl{axes.scatter()} functions.

\begin{enumerate}
	\item Load the data into a single NumPy array, using \pyinl{numpy.loadtxt()}. Mathematically, we will call this data array $\mtrx{X}$. \textit{sepal length}, \textit{sepal width}, \textit{petal length} and \textit{petal width}.
	\item Center and normalize the data by subtracting from each \textbf{column} of $\mtrx{x}$ its mean, then dividing it by its standard deviation. Use the functions \pyinl{numpy.mean()} and \pyinl{numpy.std()}, respectively. Which axis should you use for these operations? Call the resulting matrix $\mtrx{B}$.
	\item Apply the SVD method to the matrix $\mtrx{B}$:
	      \begin{minted}{python}
U, S, VT = numpy.linalg.SVD(B, full_matrices=False)
    \end{minted}
	      (where \pyinl{U, S, VT} are, respectively, the matrices $\mtrx{U},\ \mtrx{\Sigma},\ $ and $\T{\mtrx{V}}$ in the SVD of $\mtrx{B}$)
	\item In the terms of PCA, the columns of $\mtrx{U}$ are called the \textit{Principal Components} (PC) of $\mtrx{B}$. The number of principal components $N$ is the number of variables in the data:
	      \[
		      \PC{1},\ \PC{2},\ \dots,\ \PC{N}.
	      \]

	      The corresponding Singular Values are the diagonal elements of the matrix $\mtrx{\Sigma}$, denoted as
	      \[
		      \sigma_{1},\ \sigma_{2},\ \dots,\ \sigma_{N}.
	      \]

	      The singular values are ordered by absolute value in descending order. The variance in $\PC{i}$ is given by
	      \[
		      \var{i} = \frac{\sigma_{i}^{2}}{r-1},
	      \]
	      where $r$ is the number of \textit{samples} and $i=1,2,\dots,N$.

	      Plot the variance for each PC in descending order.
	\item What do we get if we divide each $\var{i}$ by the sum of all variance terms? Plot these values as function of $i$. (\textbf{hint}: what is the sum of these terms?)
	\item Plot the \textit{cumulative sum} of the values calculated in the previous step, again against $i$ (use the function \pyinl{numpy.cumsum()}). Does the cumulative sum seem to converge, and if so - to what limit? Explain briefly.
	\item Since the PCs are prdered by their variance in descending order, we can choose the first two PCs, $\PC{1}$ and $\PC{2}$, and project the data on to them. This will keep the most variance in the data while allowing us to visualize it. The matrix $\mtrx{U}$ has the PCs as its columns, and $\mtrx{\Sigma}$ are their respective values - both are sorted by variance in descending order. Therefore, the full data in terms of the PCs is given by
    \begin{equation}
        \mtrx{P} = \mtrx{U} \cdot \mtrx{\Sigma}.
        \label{eq:principal_components_US}
    \end{equation}
    The projection of the data on the first $k$ PCs is then given by the matrix product of the frst $k$ rows of $\mtrx{U}$ and the first $k$ rows of the matrix $\mtrx{\Sigma}$:
    \begin{equation}
        \mtrx{P}_{j=1}^{k} = \mtrx{U}_{j=1}^{k} \cdot \mtrx{\Sigma}_{j=1}^{k}.
        \label{eq:principal_components_first_k_dims}
    \end{equation}
    Implement this last projection (\autoref{eq:principal_components_first_k_dims}) on the first two PCs (i.e. $k=2$), and plot the result as a scatter plot of $\PC{2}$ vs. $\PC{1}$. Use \pyinl{axes.scatter(x, y)}, where \pyinl{x} and \pyinl{y} are the components of the data points in the $\PC{1}$ and $\PC{2}$ basis verctors, respectively.
    \item You can color the data using the last column, which separates the samples according to the plant's type: Setosa, Versicolor and Virginica. These are represented in the data set as \pyinl{0}, \pyinl{1} and \pyinl{2}, respectively. To use these columns as the color in the scatter plot from the previous step, use the following steps:
        \begin{enumerate}
            \item Import the following two functions from MatPlotLib's \pyinl{colors} module:
            \begin{minted}{python}
from matplotlib.colors import ListedColormap, BoundaryNorm
            \end{minted}

        \item After loading the data into a matrix, save the last column as the plant's type (make sure this array contains only the values \pyinl{0}, \pyinl{1} and \pyinl{2}).
            \begin{minted}{python}
species = data[:, -1]
            \end{minted}

        \item Generate a color map where \pyinl{0} is mapped to \pyinl{red}, \pyinl{1} to \pyinl{green} and \pyinl{2} to \pyinl{blue} (or whatever other three colors you want):
            \begin{minted}{python}
colors = ["red", "green", "blue"]
cmap = ListedColormap(colors)
bounds = [-0.1, 0.9, 1.9, 2.9]
norm = BoundaryNorm(bounds, cmap.N)
            \end{minted}

        \item Apply the color map to your scatter plot (note that the first two arguments are hidden since they relate to the task in th previous step):
            \begin{minted}{python}
ax.scatter(..., ..., c=species, cmap=cmap, norm=norm)
            \end{minted}
    \end{enumerate}

\end{enumerate}
\end{document}
